\documentclass[12pt,a4paper,twoside]{report}
\usepackage[utf8]{inputenc}
\usepackage[T1]{fontenc}
\usepackage[english]{babel}
\usepackage{graphicx}
\usepackage{amsmath, amssymb, amsthm}
\usepackage{cite}
\usepackage[colorlinks=true, linkcolor=black, citecolor=black, urlcolor=black]{hyperref}
\usepackage[linesnumbered,ruled,vlined]{algorithm2e}
\usepackage[nameinlink]{cleveref}
\usepackage{xspace}
\usepackage[colorinlistoftodos]{todonotes}
\usepackage[scaled=.83]{beramono} 
\usepackage{lineno}
\usepackage{tikz-inet}
\usepackage{stmaryrd}
\usepackage{float}
\usepackage{tablefootnote}

\usetikzlibrary{calc}
\linenumbers
\crefname{algocf}{alg.}{algs.}
\Crefname{algocf}{Algorithm}{Algorithms}
\crefformat{footnote}{#2\footnotemark[#1]#3}

% Linear
\newcommand{\agentLinear}[4][]{% [#1=options, #2=name, #3=q, #4=r]
	\inetcell[arity=1, #1](#2){$\mathit{L}_{#3,#4}$}
}
% Concrete
\newcommand{\agentConcrete}[3][]{% [#1=options, #2=name, #3=k]
	\inetcell[arity=0, #1](#2){$\mathit{C}_{#3}$}
}
% Add
\newcommand{\agentAdd}[2][]{% [#1=options, #2=name]
	\inetcell[arity=2, #1](#2){$\mathit{A}$}
}
% Mul
\newcommand{\agentMul}[2][]{% [#1=options, #2=name]
	\inetcell[arity=2, #1](#2){$\mathit{M}$}
}
% ReLU
\newcommand{\agentReLU}[2][]{% [#1=options, #2=name]
	\inetcell[arity=1, #1](#2){$\mathit{R}$} 
}
% AddCheckLinear
\newcommand{\agentAddCheckLinear}[4][]{% [#1=options, #2=name, #3=q, #4=r]
	\inetcell[arity=2, #1](#2){$\mathit{AL}_{#3,#4}$}
}
% MulCheckLinear
\newcommand{\agentMulCheckLinear}[4][]{% [#1=options, #2=name, #3=q, #4=r
	\inetcell[arity=2, #1](#2){$\mathit{ML}_{#3,#4}$}
}
% AddCheckConcrete
\newcommand{\agentAddCheckConcrete}[3][]{% [#1=options, #2=name, #3=k]
	\inetcell[arity=1, #1](#2){$\mathit{AC}_{#3}$}
}
% MulCheckConcrete
\newcommand{\agentMulCheckConcrete}[3][]{% [#1=options, #2=name, #3=k]
	\inetcell[arity=1, #1](#2){$\mathit{MC}_{#3}$}
}
% Materialize
\newcommand{\agentMaterialize}[2][]{% [#1=options, #2=name]
	\inetcell[arity=1, #1](#2){$\mathit{MAT}$}
}
% Eraser
\newcommand{\agentEraser}[2][]{% [#1=options, #2=name]
	\inetcell[arity=0, #1](#2){$\epsilon$}
}
% Duplicator
\newcommand{\agentDuplicator}[2][]{% [#1=options, #2=name]
	\inetcell[arity=2, #1](#2){$\delta$}
}
% TermAdd
\newcommand{\agentTermAdd}[2][]{% [#1=options, #2=name]
	\inetcell[arity=2, #1](#2){$\mathit{TA}$}
}
% TermMul
\newcommand{\agentTermMul}[2][]{% [#1=options, #2=name]
	\inetcell[arity=2, #1](#2){$\mathit{TM}$}
}
% TermReLU
\newcommand{\agentTermReLU}[2][]{% [#1=options, #2=name]
	\inetcell[arity=1, #1](#2){$\mathit{TR}$}
}
% Symbolic
\newcommand{\agentSymbolic}[3][]{% [#1=options, #2=name, #3=symbol]
	\inetcell[arity=0, #1](#2){$\mathit{S_{#3}}$}
}
% Interaction Rule
\newcommand{\interactionrule}[3][]{
	\begin{tikzpicture}
		\node(arrow) at (0,0){
			\begin{tikzpicture}
				\node (arrow) at (0,0) {$\longrightarrow$};
				\if\relax\detokenize{#1}\relax\else
					\node[above=0.1cm of arrow, font=\scriptsize\itshape] {#1};
				\fi
			\end{tikzpicture}
		};
		% LHS
		\node(LHS)[left=of arrow]{
			\begin{tikzpicture} #2 \end{tikzpicture}
		};
		% RHS
		\node(RHS)[right=of arrow]{
			\begin{tikzpicture} #3 \end{tikzpicture}
		};
	\end{tikzpicture}
}

\input{cmds}

\title{VEIN: VErification via Interaction Nets for Neural Networks}
\author{Eric Marin}
\date{\today}

\begin{document}

\maketitle

\begin{abstract}
	This thesis introduces \textbf{VEIN} (VErification via Interaction Nets), a framework for neural network
	verification with a focus on neural network equivalence. It acts as a	formally verified
	preprocessor that reduces neural networks to a normal form before they can be compared by a solver.
	To enable this reduction, \textbf{VEIN} translates neural networks into Interaction Nets, a graph rewriting
	computational model, and applies a set of graph rewriting rules to reduce the network into an
	Abstract Syntax Tree. We present the complete framework and provide in detail: the translation
	process, the graph rewriting rules and rigorous proofs of both soundness and termination for the
	reduction process.
\end{abstract}

\tableofcontents

% Introduction: Context -> Problem -> Solution -> Validation -> Outline
\chapter{Introduction}
\label{ch:introduction}

This chapter provides an overview of the thesis...

% This gives the reader useful information to understand the problem
Verification of neural networks is very important...

% This tells the reader the problem your solution addresses
However, larger networks can be hard to process...

% This tells the reader what you did
We show a simple method to simplify big networks...

% This tells the reader why what you did is correct
We show that our simplification is sound...

% This tells the reader the organisation of the thesis, namely what sections are there and what they contain
\section{Outline}
\label{sec:outline}
This section outlines the organization of the thesis:
\begin{itemize}
	\item \textbf{\Cref{ch:background}}: Introduces key concepts to understand the contributions...
	\begin{itemize}
		\item \textbf{\Cref{sec:neural-networks}}: Introduces the concepts of neural networks and deep learning...
		\item \textbf{\Cref{sec:interaction-nets}}: Introduces the interaction net computational model...
		\item \textbf{\Cref{sec:satisfiability-modulo-theories}}: Introduces SMT solvers...
	\end{itemize}
	\item \textbf{\Cref{ch:core}}: Dives deeper into the implementation...
	\item \textbf{\Cref{ch:related-work}}: Analyzes existing verification approaches...
	\item \textbf{\Cref{ch:conclusion}}: Summarizes the findings and discusses future work...
\end{itemize}


% Background
\chapter{Background}
\label{ch:background}

% This chapter covers the fundamental concepts and notions required to understand the work

\section{Neural Networks}
\label{sec:neural-networks}

% This section introduces Neural Networks


\section{Interaction Nets}
\label{sec:interaction-nets}

% This section introduces Interaction Nets


\section{Satisfiability Modulo Theories}
\label{sec:satisfiability-modulo-theories}

% This section introduces Satisfiability Modulo Theories



% Core
\chapter{Core}
\label{ch:core}

% This chapter contains the core of the thesis, where the work is presented in details.


% Related Work
\chapter{Related Work}
\label{ch:related-work}

% This chapter discusses related literature and positions this work within the field.


% Conclusion
\chapter{Conclusion}
\label{ch:conclusion}

% This chapter provides concluding remarks and potentially directions for future work.


\bibliographystyle{plain}
\bibliography{references}

\end{document}
